\section{Третья программа}
\subsection{Небуферизованный ввод-вывод}
\subsubsection{Исходный код}
Исходный код третьей программы с небуферизованным вводом-выводом 
представлен в листинге~\ref{lst:prog-3-1}.

\begin{listing}[Исходный код третьей программы с небуферизованным вводом-выводом]{lst:prog-3-1}{C++}
#include <fcntl.h>
#include <unistd.h>
#include <sys/stat.h>
#include <stdio.h>

void print_stat(char *msg) 
{
    struct stat statbuf;
    stat("q.txt", &statbuf);
    printf("%s: inode = %lu, size = %ld bytes\n", msg, statbuf.st_ino, statbuf.st_size);
}
int main() 
{
    int fd1 = open("q.txt", O_RDWR);
    print_stat("open fd1");
    int fd2 = open("q.txt", O_RDWR);
    print_stat("open fd2");
    for (char c = 'a'; c <= 'z'; c++) 
    {
        if (c % 2)
            write(fd1, &c, 1);
        else
            write(fd2, &c, 1);
        print_stat("write");
    }
    close(fd1);
    print_stat("close fd1");
    close(fd2);
    print_stat("close fd2");
    return 0;
}
\end{listing}

\subsubsection{Результат работы}
Результат работы третьей программы с небуферизованным вводом-выводом
представлен на рисунке~\ref{img:3-1.png}.
\img{15cm}{3-1.png}{Результат работы второй программы с небуферизованным вводом-выводом}

\subsubsection{Анализ полученного результата}
% Системный вызов open создаёт два дескриптора открытого файла в режиме 
% <<чтения и записи>> (fd1, fd2). Соответствующие дескрипторам структуры 
% file ссылаются на один и тот же inode файла <<q.txt>>. Ссылки на эти 
% структуры помещаются в таблицу открытых процессом файлов под индексами 
% 3 и 4, так как дескрипторы 0, 1 и 2 заняты стандартными потоками.

% В первой итерации цикла системный вызов write для дескриптора fd1 записывает 
% символ <<a>> в открытый файл, что приводит к увеличению на 1 значения поля 
% f\_pos в соответствующей структуре file. Размер файла становится равным 1 байту.

% Так как дескрипторам fd1 и fd2 соответствуют независимые структуры file, то 
% для fd2 значение поля f\_pos остаётся равным 0. Поэтому на второй итерации 
% цикла вызов write для fd2 заменит символ <<a>> в файле на символ <<b>>. 
% Размер файла остаётся равным 1 байту, а значение f\_pos для fd2 увеличится на 1.
% Поочерёдная запись символов в файл продолжается до конца цикла.

% Таким образом, в файл будет записана только каждая вторая буква алфавита, 
% начиная с <<b>>, и итоговый размер файла составит 13 байт, так как вызов 
% write для fd2 будет заменять символы, записанные вызовом write для fd1.

В данном варианте программы один и тот же целевой файл открывается дважды системным 
вызовом open() в режиме чтения и записи (O\_RDWR). В результате этого в таблице открытых 
файлов процесса выделяются два уникальных целочисленных дескриптора (fd1 и fd2). 
На уровне ядра операционной системы создаются две независимые структуры struct file, 
обе из которых ссылаются на один и тот же inode открываемого файла. Важнейшим аспектом 
является то, что каждая из этих структур обладает собственным, независимым указателем 
смещения f\_pos, который инициализируется нулем.

В цикле программа поочередно выполняет системные вызовы write() для каждого из дескрипторов. 
При первой итерации вызов write() через дескриптор fd1 записывает первый символ по нулевому 
смещению. После выполнения операции f\_pos первой структуры struct file инкрементируется на 
1 байт. Текущий размер файла на диске становится равным 1 байту.

Однако при следующей итерации системный вызов write() выполняется уже для дескриптора fd2. 
Поскольку его ассоциированная структура struct file независима, ее указатель f\_pos 
всё еще равен нулю. Следовательно, запись нового символа происходит с нулевой позиции, 
физически перезаписывая ранее сохраненный символ. После этой операции f\_pos 
второй структуры также увеличивается на 1, а размер файла остается равным 1 байту.

Такая конкурентная поочередная запись продолжается до конца цикла. Из-за наличия 
двух независимых, несинхронизированных указателей файлового смещения, данные, 
передаваемые ядру через fd2, систематически перезаписывают информацию, только что 
записанную через fd1. В результате происходит потеря данных: итоговый файл имеет 
размер всего 13 байт и содержит только те символы, которые были выведены через второй 
файловый дескриптор, то есть каждую букву алфавита через одну позицию, начиная с <<b>>.

Во избежание потерь данных при записи необходимо использовать флаг 
O\_APPEND при использовании функции open.

\subsection{Буферизованный ввод-вывод}
\subsubsection{Исходный код}
Исходный код второй программы с буферизованным вводом-выводом 
представлен в листинге~\ref{lst:prog-3-2}.

\begin{listing}[Исходный код второй программы с буферизованным вводом-выводом]{lst:prog-3-2}{C++}
#include <fcntl.h>
#include <stdio.h>
#include <sys/stat.h>

void print_stat(char *msg) 
{
    struct stat statbuf;
    stat("q.txt", &statbuf);
    printf("%s: inode = %ld, size = %ld bytes\n", msg, statbuf.st_ino, statbuf.st_size);
}
int main() 
{
    FILE *fs1 = fopen("q.txt", "w");
    print_stat("fopen fs1");
    FILE *fs2 = fopen("q.txt", "w");
    print_stat("fopen fs2");
    for (char c = 'a'; c <= 'z'; c++) 
    {
        if (c % 2)
            fprintf(fs1, "%c", c);
        else
            fprintf(fs2, "%c", c);
        print_stat("fprintf");
    }
    fclose(fs1);
    print_stat("fclose fs1");
    fclose(fs2);
    print_stat("fclose fs2");
    return 0;
}
\end{listing}

\subsubsection{Результат работы}
Результат работы второй программы с буферизованным вводом-выводом 
представлен на рисунке~\ref{img:3-2.png}.
\img{15cm}{3-2.png}{Результат работы второй программы с буферизованным вводом-выводом}

\subsubsection{Анализ полученного результата}
% Вызов библиотечной функции fopen откроет файл <<q.txt>> в режиме 
% <<только для записи>> и вернёт два указателя на инициализированные 
% структуры FILE (fs1, fs2). Во время вызова создаются две незваисимые 
% структуры file с дескрипторами 3 и 4, значения которых будут записаны 
% в поля \_fileno соотвествующих структур \_IO\_FILE. Эти структуры 
% file ссылаются на один и тот же inode файла <<q.txt>>.

% Вызов библиотечной функции fprintf для fs1 записывает символ <<a>> 
% в буфер, что приводит к смещению в соответствующей структуре указателя 
% \_IO\_WRITE\_PTR на 1 байт.

% Вызов библиотечной функции fprintf для fs2 записывает символ <<b>> в 
% буфер, что приводит к смещению в соответствующей структуре указателя 
% \_IO\_WRITE\_PTR на 1 байт.

% В результате работы цикла в буфере fs1 будет находиться каждая вторая 
% буква алфавита, начиная с <<a>>, а в fs2 --- каждая вторая буква 
% алфавита, начиная с <<b>>.

% Вызов fclose для fs2 приводит к записи из буфера в открытый файл. 
% Значение поля f\_pos для дескриптора 4 и размера файла становится 13. 
% После этого дескриптор закрывается.

% Так как были созданы две структуры открытых файлов, ссылающихся на 
% один и тот же inode, а в структуре, соответствующей fs1, значение 
% поля f\_pos осталось равным 0, то при вызове fclose для fs1 буфер 
% записывается в открытый файл с начала, что приводит к замене 
% последовательности символов буфера fs2 на последовательность 
% символов буфера fs1. Размер файла остаётся неизменным (13 байт).

% Таким образом, итоговое содержимое файла зависит от порядка вызовов 
% fclose. В данном случае в файл будет записан каждый второй символ 
% алфавита, начиная с буквы <<a>>, а размер файла составит 13 байт.

В этой версии программы целевой файл дважды открывается библиотечной функцией 
fopen() в режиме <<только для записи>>. Аналогично предыдущему случаю, ядро 
операционной системы создает два файловых дескриптора и две независимые структуры 
struct file с начальными смещениями f\_pos = 0. Фундаментальное отличие заключается 
в том, что библиотека стандартного ввода-вывода создает для каждого дескриптора 
структуру FILE (fs1 и fs2), выделяя для них собственные локальные буферы в адресном 
пространстве пользователя.

Внутри цикла вызовы библиотечной функции fprintf() не инициируют немедленных обращений 
к ядру. Вместо этого происходит накопление данных в пользовательских буферах. При 
обращении к fs1 символ помещается в первый буфер (сдвигая внутренний указатель 
\_IO\_WRITE\_PTR), а при обращении к fs2 --- во второй. Таким образом, по завершении 
цикла в оперативной памяти процесса формируются два независимых массива данных по 
13 байт каждый, при этом реальный файл на диске остается пустым.

Сброс накопленной информации из адресного пространства пользователя на физический 
носитель происходит только при закрытии потоков функцией fclose().
Когда выполняется вызов fclose(fs1), библиотека инициирует системный вызов write(), 
который единым блоком записывает 13 байт из первого буфера в файл. Запись 
осуществляется с позиции f\_pos = 0 первой структуры struct file, после чего её 
смещение сдвигается на 13 байт.

Следом выполняется вызов fclose(fs2). Он также инициирует системный вызов write() 
для сброса 13 байт из второго буфера. Однако, поскольку вторая структура struct 
file абсолютно независима и ее смещение f\_pos всё еще равно нулю, ядро записывает 
эти 13 байт с самого начала файла. Это приводит к полной перезаписи всего блока 
данных, только что сохраненного вызовом fclose(fs1). В итоге на выходе формируется 
файл размером ровно 13 байт, содержащий исключительно последовательность символов 
из буфера, ассоциированного с fs2.

Во избежание потерь данных при записи необходимо использовать флаг <<a>>
при использовании функции fopen. 

\subsection{Схема связи структур}
\img{8cm}{3.pdf}{Схема связи структур в первой программе}
